%; whizzy paragraph -pdf xpdf -latex ./whizzypdfptex.sh
%; whizzy-paragraph "^\\\\begin{frame}\\|\\\\emtext"
% latex beamer presentation.
% platex, latex-beamer でコンパイルすることを想定。 

%     Tokyo Debian Meeting resources
%     Copyright (C) 2012 Junichi Uekawa

%     This program is free software; you can redistribute it and/or modify
%     it under the terms of the GNU General Public License as published by
%     the Free Software Foundation; either version 2 of the License, or
%     (at your option) any later version.

%     This program is distributed in the hope that it will be useful,
%     but WITHOUT ANY WARRANTY; without even the implied warreanty of
%     MERCHANTABILITY or FITNESS FOR A PARTICULAR PURPOSE.  See the
%     GNU General Public License for more details.
%     You should have received a copy of the GNU General Public License
%     along with this program; if not, write to the Free Software
%     Foundation, Inc., 51 Franklin St, Fifth Floor, Boston, MA  02110-1301 USA

\documentclass[cjk,dvipdfmx,12pt]{beamer}
\usetheme{Tokyo}
\usepackage{monthlypresentation}

%  preview (shell-command (concat "evince " (replace-regexp-in-string "tex$" "pdf"(buffer-file-name)) "&")) 
%  presentation (shell-command (concat "xpdf -fullscreen " (replace-regexp-in-string "tex$" "pdf"(buffer-file-name)) "&"))
%  presentation (shell-command (concat "evince " (replace-regexp-in-string "tex$" "pdf"(buffer-file-name)) "&"))

%http://www.naney.org/diki/dk/hyperref.html
%日本語EUC系環境の時
\AtBeginDvi{\special{pdf:tounicode EUC-UCS2}}
%シフトJIS系環境の時
%\AtBeginDvi{\special{pdf:tounicode 90ms-RKSJ-UCS2}}

\newenvironment{commandlinesmall}%
{\VerbatimEnvironment
  \begin{Sbox}\begin{minipage}{1.0\hsize}\begin{fontsize}{8}{8} \begin{BVerbatim}}%
{\end{BVerbatim}\end{fontsize}\end{minipage}\end{Sbox}
  \setlength{\fboxsep}{8pt}
% start on a new paragraph

\vspace{6pt}% skip before
\fcolorbox{dancerdarkblue}{dancerlightblue}{\TheSbox}

\vspace{6pt}% skip after
}
%end of commandlinesmall

\title{東京エリアDebian勉強会}
\subtitle{第150回 2017年4月度}
\author{杉本 典充  dictoss@live.jp}
\date{2017年4月22日}
\logo{\includegraphics[width=8cm]{image200607/openlogo-light.eps}}

\begin{document}

\begin{frame}
\titlepage{}
\end{frame}

\begin{frame}{Agenda}
 \begin{minipage}[t]{0.45\hsize}
  \begin{itemize}
  \item 最近あったDebian関連のイベント報告
	\begin{itemize}
	\item 第148回東京エリアDebian勉強会
    \item OSC 2017 Tokyo/Spring (第149回東京エリアDebian勉強会)
	\end{itemize}
  \item 事前課題発表
  \end{itemize}
 \end{minipage}
 \begin{minipage}[t]{0.45\hsize}
  \begin{itemize}
   \item Debian Trivia Quiz
  \end{itemize}
 \end{minipage}
\end{frame}

\section{イベント報告}
\emtext{イベント報告}

\begin{frame}{第148回東京エリアDebian勉強会}
\begin{itemize}
\item 20117年2月11日に、朝日ネットさんをお借りして開催しました
\item 参加者は9名でした
\item Bug Squash partyを行いました
  \begin{itemize}
    \item \url{https://wiki.debian.org/BSP/2017/02/jp/Tokyo}
  \end{itemize}
\end{itemize} 
\end{frame}

\begin{frame}{OSC 2017 Tokyo/Spring (第149回東京エリアDebian勉強会)}
\begin{itemize}
\item 2017年3月11日に、明星大学様を会場にイベントが開催されました
\item 東京エリアDebian勉強会/DebianJPとして出展しました
\item セミナー発表は「Debian updates」という表題で、岩松さんが発表しました
\item セミナーの参加者は25名でした
\end{itemize} 
\end{frame}

\section{事前課題}
\emtext{事前課題}
{\footnotesize
 \begin{prework}{ user }
  \begin{enumerate}
  \item aaa
  \end{enumerate}
\end{prework}

}

\subsection{問題}
%; whizzy-master ../debianmeetingresume201311.tex
% $B0J>e$N@_Dj$r$7$F$$$k$?$a!"$3$N%U%!%$%k$G(B M-x whizzytex $B$9$k$H!"(Bwhizzytex$B$,MxMQ$G$-$^$9!#(B
%

\santaku
{2017$BG/$N(BDebian Project Leader(=DPL)$BA*5s$,9T$o$l!"EjI<7k2L$,H/I=$5$l$^$7$?!#(BDPL$B$KA*$P$l$?$N$OC/$G$7$g$&$+!#(B}
{Mehdi Dogguy}
{Nobuhiro Iwamatsu}
{Chris Lamb}
{C}
{Chris Lamb$B$5$s$,A*=P$5$l$^$7$?!#$*$a$G$H$&$4$6$$$^$9!*(B\url{https://bits.debian.org/2017/04/results-dpl-elections-2017.html}$B!#$J$*!"(B2017$BG/(BDPL$BA*5s$N>pJs$O<!$N(Bweb$B%Z!<%8$r;2>H$/$@$5$$(B\url{https://www.debian.org/vote/2017/vote_001}$B!#(B}


\section{セミナー1}
\emtext{Debianパッケージのカスタマイズについて}

\section{Hack time}
\emtext{Hack time}

%\begin{frame}{成果を記入下さい!}
%\end{frame}
  
\section{今後のイベント}
\emtext{今後のイベント}
\begin{frame}{今後のイベント}
  \begin{itemize}
  \item 2017/5/20(土) 東京エリアDebian勉強会（予定）
    \begin{itemize}
    \item 日時、会場、セミナー内容は調整中
    \end{itemize}
  \item 2017/6/17(土) 東京エリアDebian勉強会（予定）
    \begin{itemize}
    \item 日時、会場、セミナー内容は調整中
    \end{itemize}
  \item 2017/7/15(土) OSC 2017 Hokkaido（未定）
    \begin{itemize}
    \item 出展、またはDebianJPメンバが参加する可能性あり
    \end{itemize}
  \item 発表者は募集中です。ぜひともよろしくお願いします。
  \end{itemize}
\end{frame}

\end{document}

;;; Local Variables: ***
;;; outline-regexp: "\\([ 	]*\\\\\\(documentstyle\\|documentclass\\|emtext\\|section\\|begin{frame}\\)\\*?[ 	]*[[{]\\|[]+\\)" ***
;;; End: ***
